\subsection{Whole original projective twists under arbitrary coefficient change}

\begin{theorem}[The genuine original degree-one pole line commutes with coefficients]
\label{projective:actual-original-pole-line-coefficient-pullback}
\lean{EnvelopeProperty.ProjectiveTwistCoefficientPullback.coordinatePolePullbackIso}
\leanok
\uses{projective:actual-coordinate-hyperplane-native-module-pullback}
For ANY homomorphism of commutative coefficient rings, the actual module
pullback of the original coordinate-hyperplane $\mathcal O(1)$ is isomorphic
to the SAME original $\mathcal O(1)$ over the target coefficient ring.
No injectivity, flatness, field, reducedness, or comparison premise is supplied.
\end{theorem}
\begin{proof}
The original pole line is the actual sheaf dual of the original coordinate
ideal module. Its invertibility is already derived from genuine local
hyperplane equations. Apply the genuine evaluation-preserving whole dual
pullback isomorphism, which derives the pullback of that original dual as
the same native dual of the pulled original module. The preceding canonical
coordinate ideal-module pullback isomorphism identifies that module with the
original coordinate ideal module over the target ring. Dualize this actual
whole isomorphism in the correct contravariant direction. Its composite is
the asserted whole original pole-line isomorphism.
\end{proof}

\begin{theorem}[Every original nonnegative twist retains its exact degree]
\label{projective:actual-original-pole-power-coefficient-pullback}
\lean{EnvelopeProperty.ProjectiveTwistCoefficientPullback.coordinatePolePowerPullbackIso}
\leanok
\uses{projective:actual-original-pole-line-coefficient-pullback}
For every original $d\in\mathbb N$, the actual pullback of the original
$\mathcal O(d)$ along any coefficient-ring map is isomorphic to the SAME
original $\mathcal O(d)$ over the new ring. The exact exponent $d$ is retained.
\end{theorem}
\begin{proof}
Induct on the actual nonnegative monoidal power defining the original twist.
At degree zero, the genuine monoidal pullback unit isomorphism gives the
asserted comparison. At the successor step, the actual pullback tensorator
identifies the pullback of the original tensor product with the tensor product
of the true pullbacks. Tensor the induction isomorphism with the preceding
actual degree-one isomorphism. The result is precisely the original successor
power over the new ring. All maps are actual whole native module isomorphisms.
\end{proof}
